\section{Limitaciones, aspectos positivos y recomendaciones al proceso productivo}

Las principales limitaciones identificadas durante la práctica, se encuentra la disponibilidad de recursos y de datos. En cuanto a los recursos, uno de los principales obstáculos surgie al intentar replicar el análisis estadístico descrito en el paper \textit{Modular Deconstruction Reveals the Dynamical and Physical Building Blocks of a Locomotion Motor Program} (2015) \cite{AMB}. Este estudio utiliza \textit{toolboxes} que no están disponibles en la licencia estudiantil de MATLAB proporcionada por la universidad, por lo que una gran cantidad de tiempo fue dedicada a programar desde cero varias de las funciones utilizadas.\\

Respecto a la disponibilidad de datos, los resultados obtenidos no son representativos debido al tamaño limitado de la muestra. Además, no se dispone de información detallada sobre la metodología utilizada para la recolección de los datos originales, lo que dificulta su validación y análisis. En particular, los datos de posición espacial de las neuronas dependen de cómo se dispone el ganglio durante la medición, y si bien las neuronas de \textit{Aplysia} son de un tamaño mayor al estándar, siguen siendo lo suficientemente pequeñas para que un poco de variabilidad espacial pueda resultar en que se encuentre en la posición de otra neurona de otra grabación, lo que limita la confiabilidad del análisis realizado.\\

Estas limitaciones impidieron cumplir completamente con la actividad planteada en el plan de práctica: \textit{``Se plantearán diagramas del fenómeno a modelar, utilizando herramientas para el modelamiento tales como MATLAB, y se definirán las ecuaciones para la construcción y formulación del modelo matemático, con base en la fenomenología del proceso a modelar''}, ya que el tiempo disponible para la realización de la práctica . No obstante, se considera que sí se logró avanzar en el desarrollo de la competencia específica CE6: \textit{``Modelar y resolver problemas complejos en las distintas áreas de aplicación de la biotecnología [...] aplicando conocimientos y herramientas científicas y tecnológicas''}, particularmente en la etapa inicial de búsqueda bibliográfica y análisis estadístico de datos experimentales.\\

Como mejora concreta al proceso, se desarrolló un código que implementa la primera parte del análisis estadístico estándar, utilizando exclusivamente funciones compatibles con la licencia estudiantil de MATLAB. Este código permite a los investigadores ahorrar una cantidad significativa de tiempo en la limpieza y análisis preliminar de datos, optimizando así los recursos disponibles para avanzar en el modelamiento del sistema.\\

Un aspecto positivo destacado fue el ambiente de colaboración dentro del laboratorio. El equipo se mostró constantemente dispuesto a apoyar en las distintas etapas del proceso investigativo. Además, se llevan a cabo reuniones departamentales cada dos semanas, donde los miembros presentan brevemente artículos científicos actuales; luego, se vota por el más interesante, y en la siguiente sesión, la persona ganadora presenta el artículo en profundidad. Esta dinámica fomenta el aprendizaje colectivo y asegura que el laboratorio se mantenga actualizado respecto a los avances recientes en el campo de la neurociencia.\\


%Dentro de las principales limitaciones encontradas durante la práctica profesional se encuentran los recursos y la disponibilidad de datos. Respecto a los recursos, el primer cuello de botella en la investigación surge al momento de querer ahorrar tiempo de investigación utilizando el análisis estadístico realizado por el paper \textit{``Modular Deconstruction Reveals the Dynamical and Physical Building Blocks of a Locomotion Motor Program, 2015''} \cite{AMB}, ya que este utiliza toolboxs a las que no se tiene acceso con la licencia estudiante proporcionada por la universidad, por la que se tuvo que programar desde cero muchas de las funciones utilizadas. Respecto a la disponibilidad de los datos, se tiene que no se logran obtener resultados representativos ya que se tienen muy pocos datos para obtener resultados, además, que no se sabe con exactitud la metodología utilizada para medir estos datos y los datos de posicion que se tiene dependen de como se coloca el ganglio al momento de medir la actividad neuronal, por lo que no es un dato representativo.\\


%Por esta limitaciones, no se logra cumplir en su totalidad la actividad propuesta en el plan de práctica "se plantearán diagramas del fenómeno a modelar, utilizando herramientas para el modelamiento tales como matlab y se definen las ecuaciones para la construcción y formulación del modelo matemático, con base en la fenomenología del proceso a modelar". A pesar de esto, se considera que de todas formas se logra cumplir con la competencia específica "CE6: Modelar y resolver problemas complejos en las distintas áreas de aplicación de la biotecnología, tales como industria, biomedicina, medioambiente, biotecnología vegetal y animal, y políticas públicas asociadas a la biotecnología, aplicando conocimientos y herramientas científicas y tecnológicas.", pues se logra la primera parte necesaria para esto, es decir, la búsqueda bibliográfica y el análisis estadístico de los datos entregados.\\

%La mejora realizada al proceso fue la creación de un código, que utiliza funciones que vienen con la licencia estudiante, que realiza toda la primera parte del análisis estadístico estándar necesario. La utilizacion de este código permite a los investigadores ahorrar meses de trabajo en el análisis y limpieza de datos para poder utilizar los recursos en el modelamiento de este.\\

%Finalmente, como aspecto positivo, se tiene que las personas que trabajan en el laboratorio siempre estan dispuestas a apoyar en el proceso de investigación del otro y que se realizan reuniones de departamento bisemanales en donde todos presentan brevemente un paper actual, se vota por el más interesante y luego a la siguiente semana la persona ganadora presenta el paper a los demás, esto permite que el laboratorio este en constante actualización de las novedades de neurociencia.